% What the Accounts Settle, and What They Leave Open: Depletion as a Cost of Production, and the Missing Step from a Rate to a Horizon
% Companion paper. Typeset from companionB_standards_horizon_v10.md; the LaTeX source, the markdown mirror and the compiled PDF are
% deposited together. Compiles with tectonic; also compatible with pdflatex and xelatex.
\documentclass[11pt]{article}
\usepackage[a4paper,margin=1in]{geometry}
\usepackage{amsmath,amssymb}
\usepackage{booktabs,longtable,array,calc}
\usepackage{xcolor}
\usepackage[colorlinks=true,allcolors=blue!45!black]{hyperref}
\usepackage{etoolbox}
\AtBeginEnvironment{longtable}{\small\setlength{\emergencystretch}{6em}\hyphenpenalty=0\tolerance=9999}
\AtBeginEnvironment{quote}{\small}
\setlength{\tabcolsep}{4pt}
\providecommand{\tightlist}{\setlength{\itemsep}{0pt}\setlength{\parskip}{0pt}}
\setcounter{secnumdepth}{-1}
\setlength{\emergencystretch}{4em}
\newcommand{\tcell}[1]{\parbox[t]{0.95\linewidth}{#1}}
\usepackage[htt]{hyphenat}
\usepackage{xurl}
\tolerance=1400
\hyphenpenalty=40
\exhyphenpenalty=0
\AtBeginEnvironment{tabular}{\setlength{\emergencystretch}{6em}\hyphenpenalty=0\tolerance=9999}
\begin{document}
\title{What the Accounts Settle, and What They Leave Open: Depletion as a Cost of Production, and the Missing Step from a Rate to a Horizon}
\author{Amin Abaee\\[0.35em]
{\small Independent Researcher}\\[0.55em]
{\small\href{https://orcid.org/\allowbreak{}0000-0002-0019-1842}{ORCID: 0000-0002-0019-1842}}\\[0.3em]
{\small\href{mailto:amin\_abaee@ut.ac.ir}{amin\_abaee@ut.ac.ir}}}
\date{September 17, 2026}
\maketitle

\emph{A commentary accompanying "Typed Flux Ledgers and Depletion Arithmetic: Conservation, Componentwise Diagnostics, and the Semantics of Depletion Horizons", and the companion to the methods study on certifying a typed ledger (Abaee, 2026d). This commentary proves no new theorem. Its arithmetic is either quoted from the main text or recomputed in Section 7 from the accounts themselves, and its statements about the standards are cited. Where the main text's labels are used, they are its own.}

\textbf{Track.} Commentary / correspondence. \textbf{Numbering convention.} Assertions in this commentary are lettered B1, B2, \ldots{} and are \emph{claims about the Record}, not results, the companion carries no theorem counter.

\vspace{0.8em}

\section*{1. The question this commentary answers}\label{1-the-question-this-commentary-answers}

The accounting standards now record depletion as a cost of production. A reader of the main text will ask whether that closes the gap the main text opens. The answer is a single sentence, which this commentary defends: \textbf{the standards settle how a depletion \emph{flow} is recorded and valued. They never take the step from a flow to a \emph{horizon}. The main text's contribution lives on that step.} Everything below is either an exact statement of what the standards do, an exhibit of what they do not reach, or a note on how to cite the difference without over-claiming.

The main text already carries the technical content: the reclassification exhibit in its Section 1.2, the horizon arithmetic and its bias in Section 10.2 (Remark 33), the aggregate's inability to certify components in Section 10.1 (Proposition 30). The classification of "time to depletion" quantities in Section 6.5. What is added here is the interface: what a statistical agency that adopts the 2025 revision can and cannot say with the main text's apparatus, and what it must still decide for itself. Section 7 carries the recomputation of the aggregate overshoot date with carbon demand excluded, on the accounts themselves, with the two conventions reported side by side.

\vspace{0.8em}

\section*{2. What changed, stated exactly}\label{2-what-changed-stated-exactly}

We read the accounts for what they settle. We name the step they do not take.\textbf{B1.} In the 2025 System of National Accounts, depletion of \emph{non-produced} natural resources is recorded as a cost of production in the current accounts. In the 2008 revision the same quantity was recorded in the account of other changes in the volume of assets and liabilities. Consequently net domestic product is affected by depreciation --- the 2025 terminology for what the 2008 revision called consumption of fixed capital. \emph{and} by depletion. So are net national income and net saving.

\textbf{B2.} The 2025 revision takes this from the SEEA Central Framework (2014), which standardised the definition and recording of depletion; the framework treats depletion of natural resources like depreciation of fixed assets, on the ground that both are used up in production.

\textbf{B3.} Two consequences are worth stating because they are routinely mis-remembered. First, the change is classification, not measurement: \textbf{no extraction rate, no physical quantity and no valuation method of the 2008 accounts is replaced.} The same number moves. Second, the effect is not confined to net aggregates in every case: recording depletion as a cost may also touch gross product where natural-resource production is undertaken by government for its own final use and output is measured by the sum-of-costs approach (2025 SNA \S7.137, as cited in the OECD compilation guide; the guide is the source used here, and the paragraph number is reproduced only as a quotation from it --- see Section 6).

\textbf{B4.} Alongside the depletion entry, the 2025 revision adopts the split-asset approach: the value of a resource, and therefore its depletion charge, is attributed between the legal owner and the extractor according to economic ownership shares, with rents to the legal owner including royalties, surtaxes and permits linked to extraction. For biological resources yielding once-only products, the boundary between cultivated (produced) and non-cultivated (non-produced) assets is redrawn, regeneration of biological resources is recorded as gross fixed capital formation, and depletion is treated as a cost of production.

\textbf{B5.} The last item deserves emphasis, because it is where the two systems differ most sharply from the quantity the main text analyses. Under the 2025 revision a fishery recovering after a closure can be recorded as \emph{investment} while its stock grows, and a fishery in decline is charged a \emph{cost}. Both are descriptions of a comparison between extraction and growth. Neither is a statement about when the resource stops supporting the economy. The accounts, on their own terms, do not ask that question.

\vspace{0.8em}

\section*{3. What the standards define is a rate}\label{3-what-the-standards-define-is-a-rate}

\textbf{B6.} The definition of depletion the 2025 revision inherits is \emph{extraction in excess of the resource's growth} --- a comparison of two rates, evaluated per accounting period, in the units of the period. It carries no horizon, no barrier, no scenario, and no probability. The main text's Section 1.2 says this of the same object: "the step from a rate to a horizon is taken nowhere in them, and it is not taken here either."

That sentence is the whole relationship between the two literatures. It is worth unpacking once, because each of the three steps across it has been taken implicitly by somebody writing about depletion.

\textbf{B7 (the first step, rate to stock).} From a rate one gets a stock-level statement only by assuming the rate: a constant-drift or a declared drift class. The main text's uniform-drift bounds (its Proposition 6, and the finite-exhaustion argument around it) are exactly the conditional form of this step: given a drift bounded in a declared class, exhaustion time lies in a stated interval, and the interval's \emph{width} is part of the result. An accounts table supplies no drift class, so it cannot be read as a bound on anything.

\textbf{B8 (the second step, stock to threshold).} A stock level is not a depletion event until a barrier is named. The main text treats the barrier as a declaration, not a discovery --- its Definition 5 (scenario conditioned hitting time) and its reserve-classification discussion both turn on this. The accounts are indifferent to which barrier a country names; so is the depletion charge.

\textbf{B9 (the third step, threshold to decision).} Whether a horizon is \emph{safe} depends on a required rate of service and a control margin. That is the main text's Theorem 24 and its Proposition 41; the standards have no counterpart, and the absence is not an oversight. A national accounting system reports what happened to values and volumes within a period; it is not a solvency test.

The three steps are the main text's subject, and none of them is optional or cosmetic: the first needs an assumption, the second needs a boundary, and the third needs a purpose. The standards need none of these because they ask none of those questions.

\vspace{0.8em}

\section*{4. An exhibit of the gap: the aggregate overshoot date}\label{4-an-exhibit-of-the-gap-the-aggregate-overshoot-date}

We put the two regimes beside the standards, and we keep the predicates apart. The one applied construction in the main text that needs no dataset recomputation is worth restating here, because it shows what a rate-based standard cannot see. Write the components of an ecological footprint as demand \(d_i\) against biocapacity \(b_i\), and form

\[ \tau_{\mathrm{agg}}=365\,\frac{\sum_i b_i}{\sum_i d_i}=365\sum_i w_i r_i,\qquad w_i=\frac{d_i}{\sum_j d_j}, \qquad r_i=\frac{b_i}{d_i}. \]

Then \(\tau_{\mathrm{agg}}\) is a demand-weighted mean of the component ratios, so \(\tau_{\mathrm{agg}}\ge\tau_{\min}:=365\min_i r_i\), with equality only where every component of positive weight coincides, and the difference

\[ \Pi_\tau=\tau_{\mathrm{agg}}-\tau_{\min}=365\sum_i w_i\bigl(r_i-r_{\min}\bigr) \]

Is a compensation premium: one-signed, exactly decomposable, unbounded in the dispersion of the components (main text, Section 10.2 and Remark 33). On the world totals of the 2025 edition of the National Footprint and Biocapacity Accounts --- the edition whose series contains the year. Since the construction Lin et al. (2018) document runs to 2014 --- the 2022 aggregate ratio is \(0.584\) of biocapacity to demand. \(\tau_{\mathrm{agg}} = 213\) d. Because the carbon component carries zero biocapacity, the minimum component ratio is \(0\). \(\Pi_\tau = 213\) d. The entire aggregate date is compensation. Restricted to the five components of positive biocapacity with renormalised weights, the same arithmetic gives \(\tau_{\mathrm{agg}} = 538\) d against \(\tau_{\min} = 365\) d, a premium of \(173\) d. With the restricted premium running \(547\) d (1961), \(346\) d (1980), \(251\) d (2000) and \(173\) d (2022): it shrinks as the components converge, as the display says it must.

\textbf{B10.} Two features of that exhibit are invisible to a rate-based standard, and they are the reason the main text exists. (i) The zero row is a \emph{published convention}: the Guidebook to the National Footprint and Biocapacity Accounts states that no biocapacity figure is computed for carbon uptake, because carbon demand is charged against forest land biocapacity and a separate carbon biocapacity would double count it (Global Footprint Network, 2021, Section 9.1.2). The convention makes \(\tau_{\min}=0\) identically, so the premium equals the whole aggregate date \emph{by construction} rather than by dispersion. A standard that records flows and values has no place to register that a headline quantity is an artefact of a bookkeeping choice. (ii) A depletion charge is additive across components; an overshoot date is not, and the direction in which the aggregate is optimistic is fixed by the weighting. That is a \emph{semantic} property of the aggregate, which is why the main text calls its subject a question of semantics rather than of accuracy.

\textbf{B11.} The same construction explains a public disagreement that is not about arithmetic. An Earth Overshoot Day for 2022 was announced as 28 July --- the 209th day of that year. The 2022 row of a later edition of the accounts gives \(\tau_{\mathrm{agg}} = 213\) d. The four-day difference is vintage and nowcasting, not a dispute about the formula: the accounts are recomputed for every year of the series at each release. Anyone comparing a published overshoot date with a published component table is comparing two vintages, and the main text's practice of pinning every quoted figure to a named release (its Section 6.5, and the vintage record in the supplementary's S5) is the least costly protection available. Section 7 now measures both effects on the accounts themselves, and finds the four-day figure smaller than what the \emph{level} of aggregation does to the same date and larger than a single edition step on a closed year --- so the practice of naming a release has to name the aggregation level with it.

\vspace{0.8em}

\section*{5. Two things a standard cannot supply, and a third it should not be asked to supply}\label{5-two-things-a-standard-cannot-supply-and-a-third-it-should-not-be-ask}

\textbf{B12 (eligibility).} The 2025 revision leaves the physical boundary of an asset to the classification of the resource, and for the renewable-resource categories it admits, what counts as an asset is drawn by viability under prevailing technology and prices. The main text's Section 1.2 reads that clause correctly: an eligibility statement of exactly the kind its Remark 36 treats as reclassification rather than as a property of the stock. A statistical standard must classify, and classification criteria are chosen, not discovered. It is a boundary on what an adopted standard can certify: a resource's presence on an asset boundary is a statement about the accounts, and its absence is not a statement about the resource.

\textbf{B13 (components).} An aggregate cannot certify its components. The main text proves this in a form any auditor can use --- its Proposition 30 computes, by linear programme, the worst deficit \(\delta_j^{\ast}(z)\) consistent with a published aggregate \(z\) and the published component bounds --- and its Section 10.1 states the consequence: an aggregate can refute or alarm, never clear. A depletion charge recorded in a country's accounts is such an aggregate. It says the resource was used up in producing this year's net product. It does not say which component of the resource base carried the shortfall, and no revision of the standard will make it say so, because the information is not in the aggregate.

\textbf{B14 (a horizon, not asked for).} It would be a category error to ask the accounts for a horizon, and the main text is explicit that the comparison runs one way: the accounts supply an \emph{instance} of the phenomenon it analyses, not a premise for it. Nothing in the main text's Sections 2 to 10 depends on any claim made here about the standards, and nothing here depends on the 2025 revision. If the revision is amended, or a country implements it differently, the main text is unchanged. That is the sense in which the interface is a one-way mirror, and it is a design property, not a hedge: the main text's objects are the ledger's own.

\vspace{0.8em}

\section*{6. Citation hygiene for anyone writing this comparison}\label{6-citation-hygiene-for-anyone-writing-this-comparison}

Three practical rules, each learned from a specific difficulty in this review cycle.

\textbf{B15. Cite no paragraph numbers of the 2025 revision.} Compilation-stage documents disagree about the structure of the text they describe: the OECD guide as published discusses the SEEA in Chapters 2, 35 and 36, while the same statement in the circulation-stage material refers to Chapters 2, 34 and 35. Paragraph references are therefore unstable at exactly the moment commentators most need them, and a wrong paragraph number in a commentary is a correction that will itself need correcting. Cite the standard by edition and by substance; where a compilation guide supplies a paragraph number, reproduce the \emph{guide} as the source of the number, as B3 does.

\textbf{B16. Distinguish endorsement from implementation.} The 2025 revision was endorsed by the United Nations Statistical Commission at its fifty-sixth session in March 2025; the schedule of country implementation, the availability of physical asset accounts, and the completeness of depletion estimates vary by country and by resource. "Adopted" and "in the published accounts of country X" are different claims, and the second generally requires the supplementary-style test of reading the numbers off the source rather than citing its title.

\textbf{B17. Never quote a depletion number as a date.} This is the main text's own negative content, and it applies with full force to the new accounts: a depletion charge is a rate, and rate times a year is a value, not a year. If a national statistical office now writes "depletion of forests, 1.4\% of GDP", that sentence is complete. The inference "so we have N years" is not in it.

For example, one further instance belongs to this section, because it is the kind of error a revision cycle hides: an exhibit quoting a 2022 figure and citing Lin et al. (2018) for it mis-credits the \emph{release}, whose series ends in 2014, even though the citation is right about the \emph{construction}. The two are different claims, and only the second is checkable against the paper's own reference list. Both this commentary's Section 4 and the main text's Remark 33 carried that ambiguity until the recomputation made the end year of the cited edition a fact on the page instead of a recollection.

\vspace{0.8em}

\section*{7. The one open computation, run}\label{7-the-one-open-computation-run}

For example, we recompute the aggregate from the tables, and we show what carbon demand does to it. This section reports the restricted and unrestricted series recomputed with carbon demand excluded, to answer three questions: whether the sign of the premium's trend survives the exclusion. How far the restricted and unrestricted dates separate over the last decade, and whether the four-day difference recorded in B11 widens or closes as the accounts are read. The arithmetic is run on the National Footprint and Biocapacity Accounts --- the 2018 edition, series 1961 to 2014, with the 2017 edition (1961 to 2013) as the revision control --- on the publisher's own \texttt{World} rows, with no interpolation and no model. \S8 gives the provenance, the hashes and the command. Two conventions, one arithmetic: ALL is the published convention, all six demand components in the ratio, in which the carbon row's zero biocapacity forces \(\tau_{\min} = 0\) and the premium to the whole date. NOC is the restricted convention, carbon demand removed from numerator and denominator with the weights renormalised over the five components of positive biocapacity.

\begin{longtable}{@{}>{\raggedright\hyphenpenalty=0\exhyphenpenalty=0\arraybackslash}p{14.5mm}@{\hspace{4.0pt}}>{\raggedleft\hyphenpenalty=0\exhyphenpenalty=0\arraybackslash}p{21.7mm}@{\hspace{4.0pt}}>{\raggedleft\hyphenpenalty=0\exhyphenpenalty=0\arraybackslash}p{31.1mm}@{\hspace{4.0pt}}>{\raggedleft\hyphenpenalty=0\exhyphenpenalty=0\arraybackslash}p{32.6mm}@{\hspace{4.0pt}}>{\raggedleft\hyphenpenalty=0\exhyphenpenalty=0\arraybackslash}p{23.3mm}@{\hspace{4.0pt}}>{\raggedleft\hyphenpenalty=0\exhyphenpenalty=0\arraybackslash}p{26.7mm}@{}}
\toprule
year & carbon share of demand & \(\tau_{\mathrm{agg}}\), carbon included & \(\tau_{\mathrm{agg}}\), carbon demand excluded & gap & restricted premium \\
\midrule
1961 & 43.9\% & 498.7 d & 889.5 d & 390.8 d & 524.5 d \\
1980 & 55.4\% & 307.1 d & 688.3 d & 381.1 d & 323.3 d \\
2000 & 55.3\% & 266.2 d & 595.3 d & 329.2 d & 230.3 d \\
2005 & 58.6\% & 237.1 d & 571.9 d & 334.9 d & 206.9 d \\
2006 & 59.5\% & 231.5 d & 572.1 d & 340.6 d & 207.1 d \\
2007 & 60.3\% & 225.9 d & 568.8 d & 342.9 d & 203.8 d \\
2008 & 60.3\% & 226.8 d & 571.2 d & 344.4 d & 206.2 d \\
2009 & 60.1\% & 230.0 d & 576.5 d & 346.5 d & 211.5 d \\
2010 & 61.3\% & 219.4 d & 567.1 d & 347.7 d & 202.1 d \\
2011 & 61.6\% & 215.8 d & 561.9 d & 346.0 d & 196.9 d \\
2012 & 61.4\% & 216.1 d & 559.3 d & 343.2 d & 194.3 d \\
2013 & 60.8\% & 215.3 d & 549.6 d & 334.3 d & 184.6 d \\
2014 & 60.2\% & 216.5 d & 543.9 d & 327.4 d & 178.9 d \\
\bottomrule
\end{longtable}

\textbf{B18.} \emph{The trend's sign survives the exclusion, and its size is not an artefact of the release.} The restricted premium on the 2018 edition runs 524 d (1961), 323 d (1980), 230 d (2000), 179 d (2014); the decade means are 464 d, 358 d, 288 d, 253 d, 216 d, 191 d, and 45 of 53 year-on-year moves are declines. The main text's series --- \(547\), \(346\), \(251\), \(173\) d, read off a later edition --- has the same sign and shape and sits 4 to 8\% above this release on the shared years. So the convergence reading in B10 does not depend on keeping carbon in the ratio set, and the 4 to 8\% release difference is the same phenomenon this commentary is about.

\textbf{B19.} \emph{The two conventions part by roughly a year, and the size of the part is an identity, not a finding.} Over the last decade of the release the restricted date exceeds the published-convention date by 340.8 d on average (327.4 to 347.7 d), a ratio of 2.41 to 2.60; that ratio is exactly \(1/(1-s_{\mathrm{carbon}})\) to within 4.4e-16, with \(s_{\mathrm{carbon}}\) between 58.6\% and 61.6\%. Removing carbon demand divides the denominator by \((1-s)\) and multiplies the reported date by its reciprocal, which is a mechanical consequence of the convention --- worth stating because the two numbers are otherwise quoted as though they answered different questions.

\textbf{B20.} \emph{The four-day gap neither dissolves nor indicts the arithmetic, and the explanation on offer is incomplete.} Across an edition step on a closed year the aggregate date moves by 1.1 d on average and 1.9 d at most, so a mature-year revision does not account for a four-day figure. But reading the publisher's \texttt{World} row instead of summing the national rows moves the \emph{same} date by up to 9.2 d on this release --- larger than the gap under discussion, and not a matter of data vintage at all. The restricted date, mean 6.3 d and maximum 9.1 d across the same edition step, is a further 5.8 times more revision-sensitive than the aggregate date, because the excluded denominator is the smaller quantity and an absolute revision is divided by less. The direction, then: for the aggregate date a few-day gap is bigger than one edition step and smaller than the choice of aggregation level; for the restricted date it is routine.

\textbf{B21.} \emph{One structural fact changes how the display in B10 reads at a world aggregate.} There, cropland and built-up land have \(b_i = d_i\) identically --- 54 of 54 years, exactly --- because the world footprint of a component whose demand \emph{is} the area is priced at world-average yields, so those two components coincide by construction; consequently \(\tau_{\min}^{\mathrm{NOC}} = 365\) d in 54 of 54 years. The restricted premium is then not a measured dispersion among components but the distance of the aggregate date past the year boundary, and the "components converge" reading has to name the level at which it was computed. The main text carries this as its Remark 37; the commentary records it because it is exactly the kind of thing a rate-based standard has no place to register (B10).

\vspace{0.8em}

\vspace{0.8em}

\section*{8. Provenance, and how to re-run it}\label{8-provenance-and-how-to-re-run-it}

The release is named because the numbers are only as durable as that naming.

\begin{itemize}\setlength{\itemsep}{2pt}
\item \textbf{Data.} The 2018 edition of the National Footprint and Biocapacity Accounts (series 1961 to 2014), the publisher's own deposit, and the 2017 edition (1961 to 2013), its earlier deposit; both under CC BY-SA 4.0. The sha256 of the tables as read: 60968f7c9959537f\allowbreak{}8e67f915aca42596\allowbreak{}62b5cd42c3a0ec02\allowbreak{}d094677b4c280ef6 and 0dd766d975cd5e85\allowbreak{}f2a2d39cff1f914b\allowbreak{}92c514186ce507cb\allowbreak{}1f721a63be57b6a5. - \textbf{Retrieval.} The datasets' public download endpoint on a data-marketplace host, no account required for those two deposits. The \textbf{current} edition is not reachable that way: the publisher distributes it free but through a registration form with an emailed link, its open data endpoints answer with authorisation errors or have moved, and the only edition mirrored on a public data portal is the 2018 one. So this is the newest openly fetchable, component-level release, not the current release, and that distinction is stated wherever the numbers are used. - \textbf{What was examined and rejected.} A single-year footprint table on a data-hosting site, set aside because it carries no year column and no pinnable edition, so it cannot support a claim about a release at all. A package-manager loader whose name matches the subject turns out to be about something else entirely. Neither was used, and both are named. - \textbf{The program.} \texttt{analysis/\allowbreak{}nfa\_tau/\allowbreak{}recompute\_ta\allowbreak{}u.\allowbreak{}py}: standard library only, no network call, no interpolation. It writes one CSV per edition plus a revision ledger and a complete stdout, and it prints the aggregation diagnostic that B20 turns on. Re-running on a newer edition is two file names, not a change of arithmetic. The directory holds the two tables as fetched, a manifest with the retrieval commands and the hashes of the archive and its extracted member, and a licence note; a re-run reproduces its three outputs byte for byte from the copies archived in the record; those copies are held there for re-running and are not redistributed with the deposit, which carries the manifest, the checksums and the two retrieval commands instead. - \textbf{Provenance of B's own Section 4.} The 2022 row quoted there (\(0.584\), \(213\) d) cannot come from Lin et al. (2018), whose release's series ends in 2014. The main text's Remark 33 names the edition the year comes from and keeps that citation for the construction it documents. This is the specific reason for the release-naming practice of B11.
\end{itemize}

\vspace{0.8em}

\vspace{0.8em}

\section*{9. What a reporter should publish}\label{9-what-a-reporter-should-publish}

We close with what the record cannot carry, because that is the part a reader must know. Four sentences of practice, which is the only place a commentary can be usefully prescriptive.

1. Publish the component table with the rate comparison, not a horizon: extraction and growth, per component, with the period named. The reader who wants a date can then build one and will know what they assumed. 2. If a date is published, publish the direction of its optimism. For any demand-weighted aggregate of component ratios, the aggregate is greater than or equal to the worst component's date, and the gap is computable from the same table (B10). One extra column is all it takes. 3. Name the release, the pull and the aggregation level, not the publisher. "Per the 2024 release, on the publisher's world aggregate" is checkable; "per national accounts" is not, once revisions are routine --- and B21 measures the level's contribution in days, so it belongs in the same sentence as the release. 4. Say which classification decisions the figure depends on. A depletion charge depends on the asset boundary, on the ownership split, and on which component carries a zero by convention. All three are bookkeeping. Naming them costs nothing and makes the figure durable across revisions.

None of the four requires the main text's machinery. All four require the distinction the main text draws between a quantity that answers one question and a quantity that is read as answering another.

\vspace{0.8em}

\vspace{0.8em}

\section*{References}\label{references}

Abaee, A, 2026. Typed Flux Ledgers and Depletion Arithmetic: Conservation, Componentwise Diagnostics, and the Semantics of Depletion Horizons. Zenodo. \url{https://doi.org/10.5281/zenodo.22554177}.

Abaee, A, 2026d. Certifying a Typed Ledger: The Predicates, the Programmes, the Vintages, and the Reproduction Bundle. Companion methods study, submitted with this commentary.

Global Footprint Network, 2021. Working Guidebook to the National Footprint and Biocapacity Accounts, 2021 edition. Global Footprint Network, Oakland, CA.

Lin, D, Hanscom, L, Murthy, A, Galli, A, Evans, M, Neill, E, Mancini, M.S, Martindill, J, Medouar, F.-Z, Huang, S, Wackernagel, M, 2018. Ecological footprint accounting for countries: Updates and results of the National Footprint Accounts, 2012--2018. Resources 7, 58. \url{https://doi.org/10.3390/resources7030058}

OECD, 2025. Measuring Natural Resources in the National Accounts: Compilation guidance for the implementation of the 2025 SNA. OECD, Paris.

United Nations, 2014. SEEA Central Framework: 2012 Technical Implementation. Statistical Papers, Series M No. 96. United Nations, New York. (Issued jointly with the European Commission, the International Monetary Fund, the Organisation for Economic Co-operation and Development and the World Bank.)

United Nations, 2025. System of National Accounts 2025. United Nations Statistics Division, New York. Endorsed by the United Nations Statistical Commission at its fifty-sixth session, March 2025. \url{https://unstats.un.org/unsd/nationalaccount/sna2025.asp}

\vspace{0.8em}

\section*{Declarations}\label{declarations}

\textbf{Funding.} None declared.

\textbf{Competing interest.} The author declares no competing interest.

\textbf{Relationship to the main text.} This commentary was written as part of the same revision cycle; its standards statements were checked against the sources above and against the OECD compilation guide, and its arithmetic is quoted from the main text, Section 10.2, without recomputation. It cites no paragraph number of the 2025 revision except where the guide itself is the source of the number.

\textbf{Use of AI assistance.} Drafting and iterative review were assisted by GLM (Z.ai), Qwen (Alibaba Cloud) and DeepSeek AI, as in the main text. All statements about the standards are the author's to confirm before submission, for the reason given in B15.
\end{document}
